%%%%%%%%%%%%%%%%%%%%%%%%%%%%%%%%%%%%%%%%%
% Diaz Essay
% LaTeX Template
% Version 2.0 (13/1/19)
%
% This template originates from:
% http://www.LaTeXTemplates.com
%
% Authors:
% Vel (vel@LaTeXTemplates.com)
% Nicolas Diaz (nsdiaz@uc.cl)
%
% License:
% CC BY-NC-SA 3.0 (http://creativecommons.org/licenses/by-nc-sa/3.0/)
%
%%%%%%%%%%%%%%%%%%%%%%%%%%%%%%%%%%%%%%%%%

%----------------------------------------------------------------------------------------
%	PACKAGES AND OTHER DOCUMENT CONFIGURATIONS
%----------------------------------------------------------------------------------------

\documentclass[11pt]{diazessay} % Font size (can be 10pt, 11pt or 12pt)

%----------------------------------------------------------------------------------------
%	TITLE SECTION
%----------------------------------------------------------------------------------------

\title{\textbf{AI and reasoning} \\ {\Large\itshape What does AI reason about?}} % Title and subtitle

\author{\textbf{Jenny Vermeltfoort}\\ \textit{s3787494} \\ \textit{Leiden University} } % Author and institution

\date{\today} % Date, use \date{} for no date

%----------------------------------------------------------------------------------------

\begin{document}

\maketitle % Print the title section

%----------------------------------------------------------------------------------------
%	ABSTRACT AND KEYWORDS
%----------------------------------------------------------------------------------------

%\renewcommand{\abstractname}{Summary} % Uncomment to change the name of the abstract to something else

%----------------------------------------------------------------------------------------
%	ESSAY BODY
%----------------------------------------------------------------------------------------

\section*{Introduction}
Within this essay I make a stance on what is artificial intelligence. The Turing Test is a metric to measure the intellegence of a machine; in the first section I conclude that it is only a measurement of human psychology. But then what can replace the Turing Test? The second section introduce the human skills that are correlated to general intellegence.

%------------------------------------------------

\section*{Turing Test}
Alan Turing proposes the question whether machines can think in his paper: "Computing Machinery and Intelligence"\cite{turing}. He supplements this question with a game called 'the imitation game', now known in  popular mouth to be the Turing Test. In this imitation game the purpose of the machine is to convince an uninformed human interregator that it is human. The notion is that when a machine can answer any arbitrary question in the style of a human, it must be intelligent. McCarthy agrees with Turing's conclusion in his paper: "What is Artificial Intelligence?"\cite{mcarthy}.

Would it not be easy to distinguish the answer of a machine from the answer of a human if it were to answer in an absolute manner? Ask the machine an arithmetic question and it would answer perfectly every time. To pass the test, we must produce a machine that does not have to answer correctly, but solely tries to fool the interrogator: it answers in a slow and inaccurate, less-intelligent manner. Accomplishing this human mannerism is not hard; Dennett shows that people are easily tricked by rather dumb programs\cite{dennett}. And even though the adaptation of communication in such a sense could showcase a different type of intelligence, it is noticeable that this shifts the objective from simulating intelligence to the intentional task of tricking the interrogator. A machine that is programmed to perfectly pass the Turing Test is then the mere result of scientific research within human psychology; reducing artificial intelligence to a static product of human intelligence. This implies that the Turing Test is not a good metric that decides the intelligence of a machine, considering that when it is cheated, it only measures the intelligence of the programmer and how adequately she understands human psychology.

\section*{Problem solving}
The Turing Test is not a decisive metric to decide artificial intelligence, but it brings to light a relation between the ability to attain any general objective and intelligence. If a machine can accomplish any task without the involvement of a human, it certainly must be considered intelligent. McCarthy, in a similar framing, writes that it is important that artificial intelligence can problem-solve in the same manner as people\cite{mcarthy}. 

Human intelligence is founded upon logical reasoning and abstract categorization to solve unfamiliar problems. We use our pool of knowledge to find patterns, experiment, observe, and derive conclusions; these skills being the basis of the scientific method and resulting in the current intellectual domain. I, therefore believe it is important to produce a metric that compiles intelligence through the lens of some set of applied skills, instead of defining a metric that purely reasons about intelligence on the basis of problem solutions.

\section*{Mathmetical logic}



%------------------------------------------------

\section*{Conclusion}


%----------------------------------------------------------------------------------------
%	BIBLIOGRAPHY
%----------------------------------------------------------------------------------------

\bibliographystyle{unsrt}

\bibliography{sample.bib}

%----------------------------------------------------------------------------------------

\end{document}
\begin{table}[h] % [h] forces the table to be output where it is defined in the code (it suppresses floating)
	\caption{Example table.}
	\centering
	\begin{tabular}{l l r}
		\toprule
		\multicolumn{2}{c}{Name} \\
		\cmidrule(r){1-2}
		First Name & Last Name & Grade \\
		\midrule
		John & Doe & $7.5$ \\
		Richard & Miles & $5$ \\
		\bottomrule
	\end{tabular}
\end{table}